% This section is intended for inclusion in the consolidated Russian appendix.
% Bibliography keys used below: vdER, Shastri, Craighero, Masuda.
% All conclusions depend on the complete rectifier, not merely on a parameter lift.
\section{Какие именно именные гипотезы затрагивает выпрямление}

В этом разделе слово «опровержение» означает предъявленный полный
математический аргумент с точными сертификатами, представленный на проверку.
Внешнее рецензирование и литературный приоритет не утверждаются.

Положим $K=\mathbb Q(\sqrt{-3})$ и
\[
\gamma_T(u)=(f(u),g(u),h(u))
=(u^3-3u,u^4-4u^2,u^5-10u).
\]
В \cite[с.~1319--1320]{vdER} сформулированы три эквивалентные
гипотезы о конкретном объекте, связанном с этим вложением. Они
предполагают неполиномиальность соответственно ядра определённого LND,
факторкольца по его срезу и расширенного кольца Риса идеала кривой.
Последняя формулировка обозначена там как Asanuma's Conjecture.
Приведённый далее аргумент, применённый к удостоверенному выпрямителю
над $K$, даёт отрицание всех трёх формулировок над $\mathbb C$.

\begin{proposition}[Явные координаты связанного LND]
Пусть $k$ --- поле характеристики нуль,
$F=(P,Q,R_3)\in\operatorname{Aut}_k(\mathbb A^3)$,
$G=F^{-1}$, $F\circ\gamma=(u,0,0)$ и $0\ne b(t)\in k[t]$.
Для
\[
A=k[t,u,X_1,X_2,X_3],\qquad
D_b=b(t)\partial_u+\sum_i f_i'(u)\partial_{X_i}
\]
существует полиномиальная система координат
$A=k[t,p,q,r,s]$, в которой $D_b=\partial_s$.
\end{proposition}

\begin{proof}
Положим $V_i=f_i(u)-b(t)X_i$ и
\[
p=P(V),\qquad q=\frac{Q(V)}{b(t)},\qquad
r=\frac{R_3(V)}{b(t)},\qquad
s=\frac{u-P(V)}{b(t)}.
\]
Все числители точно делятся на $b(t)$, поскольку
$V\equiv\gamma(u)\pmod{b(t)}$. Обратная замена задаётся формулами
\[
u=p+b(t)s,\qquad
X_i=\frac{f_i(p+b(t)s)-G_i(p,b(t)q,b(t)r)}{b(t)}.
\]
Она также полиномиальна благодаря $G(p,0,0)=\gamma(p)$.
Обе композиции тождественны после обращения $b(t)$; инъективность
локализации доказывает их тождественность в исходном кольце. Наконец,
$D_b(V_i)=0$, откуда $D_b(p)=D_b(q)=D_b(r)=0$ и $D_b(s)=1$.
\end{proof}

Для модели Шастри и $b(t)=t$ получаем
\[
\ker D=K[t,p,q,r]\simeq K^{[4]}.
\]
После расширения скаляров это противоречит неполиномиальности,
утверждаемой в Conjecture~5.1 \cite{vdER}.

Исходный полином параметра Шастри равен
\[
F_0(x,y,z)=yz-x^3-5xy+2z-7x,
\]
а срез в Conjecture~5.2 равен
$\sigma=(u-F_0(V))/t$. Так как $D(\sigma)=D(s)=1$,
разность $\sigma-s$ принадлежит $B=\ker D$.
Следовательно, $A=B[\sigma]$; в частности,
\[
A/(\sigma-c)\simeq B\simeq K^{[4]}\qquad(c\in K).
\]
Это отрицание Conjecture~5.2 после перехода к $\mathbb C$.
При этом сам $F_0$ не является координатой $\mathbb A^3$:
его градиент исчезает в точках $(1,-2,5)$ и $(-1,-2,-5)$.
Некоординатность $F_0$ и координатность $\sigma$ совместимы,
поскольку речь идёт о разных полиномах в разных кольцах.

Наконец, при $R_0=K[x,y,z]$ и $I=I_{\gamma_T}$ имеем
$R_0=K[P,Q,R_3]$, $I=(Q,R_3)$, поэтому
\[
R_0[t,I/t]=K[t,P,Q/t,R_3/t]\simeq K^{[4]}.
\]
Здесь $R_0[t,I/t]\subset R_0[t,t^{-1}]$ --- именно расширенное
кольцо Риса, фигурирующее в Conjecture~5.3. Это отрицание третьей
формулировки, а не результат о произвольных кольцах Риса.

Таким образом, устранён один конкретный кандидат в трёх связанных
представлениях. Общая проблема сокращения из этого не решается.
Не решается и общая задача выпрямления $\mathbb A^1\hookrightarrow
\mathbb A^3$ либо общая гипотеза существования невыпрямимых
пространственных прямых.

\paragraph{Исторические различия.}
Предположение о комплексной невыпрямимости именно $\gamma_T$
приписано Шастри в \cite[с.~1319]{vdER}; предъявленный выпрямитель
его опровергает. Но вопросы Q.1 и Q.2 самой статьи
\cite[с.~12]{Shastri} уже решены Шастри: они относятся к
полиномиальным представлениям вещественных узлов, комплексирующимся
в замкнутые вложения, а не к их комплексной выпрямимости.
Специальные гипотезы Абъянкара, сформулированные через степени,
имели прежние опровержения, описанные в
\cite[введение, с.~115--116]{Craighero}.

\paragraph{Поля и действия.}
Полученные координаты определены над $K$ и затем над $\mathbb C$;
спуск на $\mathbb Q$ или $\mathbb R$ не доказан.
При $b=t$ естественное действие
$t\mapsto\lambda t$, $u\mapsto u$, $X_i\mapsto\lambda^{-1}X_i$
линейно в координатах ядра с весами $(1,0,-1,-1)$.
Это конкретная линеаризация. Утверждения о всех групповых действиях
из неё не следуют. В обратных связях с вложением существенна
градуировка, как и в работе \cite{Masuda}; её нельзя заменять
лишь абстрактной полиномиальностью ядра или равенством $\operatorname{ML}(B)=k$.

\subsection{Полный ML, относительный ML и фильтрация}
Для выпрямленной модели $B=k[t,p,q,r]$ четыре координатных
дифференцирования дают
\[
\ML(B)=k,\qquad \ML^*(B)=k,\qquad \operatorname{Dk}(B)=B.
\]
Здесь $\ML^*$ использует только LND со срезами, а $\operatorname{Dk}$
есть алгебра, порождённая ядрами всех ненулевых LND.
Для $D_b=\partial_s$ имеем $\pl(D_b)=B$,
$\ker D_b^{n+1}=\bigoplus_{j=0}^n Bs^j$ и
$D_b^n(\ker D_b^{n+1})=B$.
Это согласуется с большой минимальной степенью среза в исходных
координатах: речь идёт о разных характеристиках.

Для произвольного вложения, когда выпрямитель неизвестен, всё же верна
дихотомия
\[
\ML(B_b)\in\{k,k[t]\},\qquad B_b=\ker D_b.
\]
Действительно, $\partial_{X_i}$ коммутируют с $D_b$ и ограничиваются
на $B_b$. Пересечение их ядер там равно
$B_b\cap k[t,u]=k[t]$. Если непостоянный $a(t)$ лежит в полном ML,
то для любого LND $\Delta$ на $B_b$ равенство
$0=\Delta a(t)=a'(t)\Delta t$ в области характеристики нуль
влечёт $\Delta t=0$. Тогда весь $k[t]$ лежит в ML.

Поэтому $\ML(B_b)=k[t]$ было бы препятствием выпрямлению.
Однако равенство $\ML(B_b)=k$ не доказывает ни полиномиальности $B_b$,
ни выпрямимости. Если с самого начала разрешать только LND,
фиксирующие $t$, относительный ML равен $k[t]$ даже в полиномиальном
случае. Для новых ядер $H(3,7,8)$ и $H(4,5,7)$ выбор ветви дихотомии
пока не выполнен.

\subsection{Что надо сохранить для обратного перехода}
Пусть $R=k[x,y,z]$, $I=I_\gamma$,
$B_\mu=R[t,I/t^\mu]$, $\mu\ge1$, $\wt(t)=1$, $\wt(R)=0$.
Вложение выпрямимо тогда и только тогда, когда существуют однородные
координаты
\[
B_\mu=k[t,p,q,r],\qquad
\wt(t,p,q,r)=(1,0,-\mu,-\mu)
\]
с \emph{тем же выделенным элементом $t$}. Обратное направление следует
из равенств
\[
\begin{gathered}
(B_\mu)_0=R=k[p,t^\mu q,t^\mu r],\\
I=t^\mu(B_\mu)_{-\mu}=R\cap tB_\mu=(t^\mu q,t^\mu r).
\end{gathered}
\]
Ограничение $p$ на кривую является аффинным параметром; его можно
нормировать до $u$.

Эквивалентная формулировка: существует однородное LND $E$ степени $-1$
на $B_\mu$ с $E(t)=1$. Прямое направление даёт $\partial_t$.
Обратно, теорема о срезе и конечная проекция
\[
\rho(b)=\sum_{j\ge0}\frac{(-t)^j}{j!}E^j(b)
\]
дают $B_\mu=(\ker E)[t]$ и градуированный изоморфизм
$\ker E\simeq B_\mu/(t)$. Последний слой полиномиален с весами
$(0,-\mu,-\mu)$ по построению нормального базиса в разделе~\ref{sec:formal}.
Значит, поднимаются требуемые однородные координаты.

Для LND $E$, без предположения однородности, условие $E(t)=1$ ещё даёт абстрактную полиномиальность
$B_\mu$, но приведённое доказательство уже не восстанавливает вложение.
Для LND $E$ с одним лишь $E(t)\ne0$ следует только $\ML(B_\mu)=k$.
Нельзя без доказательства брать внутреннюю однородную компоненту
произвольного LND и объявлять её локально нильпотентной.

\paragraph{Непосредственная связь с плоской теоремой.}
Если идеал $I$ содержит полиномиальную координату, кривую можно
поместить в координатную плоскость и применить теорему
Абъянкара--Моха--Судзуки. Если координатой является полином параметра
$P\circ\gamma=u$, то при $R=k[P,U,V]$ нужный выпрямитель сразу равен
$(P,U-U(\gamma(P)),V-V(\gamma(P)))$.
Обычного, не обязательно координатного полинома параметра для обоих
рассуждений недостаточно.
